\documentclass[12pt]{article}

\usepackage{amsmath}
\usepackage{graphicx}
\usepackage{verbatim}

\graphicspath{/images}

\title{Assignment 4 Workout}
\author{Thomas Vroom - i6291496}
\date{October 2022}

\begin{document}
\maketitle
\pagenumbering{gobble}

\section*{Question 1}
Notes:

\begin{itemize}
    \item $P(\text{win}_{\text{normal}}) = \frac{1}{37}$
    \item $P(\text{win}_{\text{cheating}}) = \frac{1}{20}$
    \item Number of (independent) plays $n = 100$
    \item $P(\text{player } y \text{ is a cheater}) \geq 0.8$
\end{itemize}

\subsection*{A - What distribution can you use to model the number of wins in 100 games?}
Each game can either be won or lost, so this is a Bernoulli trial.
Because we are tracking the wins over multiple games, we are tracking repeated Bernoulli trials.
This corresponds to a \textbf{Binomial Distribution}.

\subsection*{B - What is the null hypothesis and its distribution?}
The null hypothesis $H_0$ is that the player is \underline{not} cheating.
This gives a Binomial distribution of $X \sim \text{Bin}(100, \frac{1}{37})$.

\subsection*{C - What is the alternative hypothesis and its \\distribution?}
The alternative hypothesis $H_A$ is that the player \underline{is} cheating.
This gives a Binomial distribution of $X \sim \text{Bin}(100, \frac{1}{20})$.

\subsection*{D - What does 80\% mean in the context of hypothesis testing?}
It means that the \textbf{power} of the test is greater than or equal to 0.8.
The power refers to the probability that we reject $H_0$ (we decide the player is a cheater) when in reality $H_A$ holds (the player is indeed a cheater):

\[ \text{power: } P(\text{reject } H_0 | H_A) \geq 0.8 \]

\subsection*{E - What is your decision rule?}
We know:

\[ P(\text{reject } H_0 | H_A) \geq 0.8 \]

\noindent
We can phrase this as: at what number of wins becomes the probability of the player being a cheater become greater than or equal to 0.8?
Or: at what $k$ becomes the sum from $k$ to 100 of $P(X = k | H_A)$ greater than or equal to 0.8:

\[ P(X \geq k | H_A) \geq 0.8 \]

\noindent
We can calculate this $k$ with the Binomial formula bin$(n, p, k)$, by summing the result of this calculation for $k = 100, 99, 98, $ etc. until the sum becomes greater than or equal to 0.8.
When this happen we simply take the $k$ where we stopped and that represents the number of games a player can win before being labeled as a cheater.
In the Binomial formula, we use $n = 100$ as the number of trials and $p = \frac{1}{20}$ as the probability of success (because $H_A$ is given in the above inequality we know we have to use cheating probability).
We can do this calculation using the following python code:

\newpage
\begin{verbatim}
from scipy.stats import binom

# number of trials
n = 100
    
# probability of winning
p = (1 / 20)
    
# power of the test
power = 0.8
    
sum = 0
k = 100
while sum < power:
    sum = sum + binom.pmf(k, n, p)
    k = k - 1
        
print("sum:", sum, "k:", str(k + 1))
\end{verbatim}

\noindent
Running this code gives $k = 3$.
This means our decision rule becomes:
If a player wins more than or exactly \textbf{3 out of a 100} games, they will be labeled as a cheater.

\subsection*{F - What is the probability of calling someone a cheater who is not cheating in the context of hypothesis testing?}
Calling someone a cheater who is not cheating is called a \textbf{Type I error}.
The probability of this happening is called the \textbf{significance level} of the test:

\[ \alpha = P(\text{reject } H_0 | H_0) \]

\newpage
\subsection*{G - What is the probability of F in your decision rule?}
We know:

\[ \alpha = P(\text{reject } H_0 | H_0) \]

\noindent
In our decision rule, we reject $H_0$ if the player wins more than or exactly 3 out of 100 games.
This means:

\[ \alpha = P(X \geq 3 | H_0) \]

\noindent
We can calculate this using the following python code:

\begin{verbatim}
from scipy.stats import binom

# number of trials
n = 100
    
# probability of winning
p = (1 / 37)
    
# starting index
i = 3
    
sum = 0
for k in range(i, n + 1):
    sum = sum + binom.pmf(k, n, p)
    
print("sum:", sum)    
\end{verbatim}

\noindent
Running this code gives us a probability of:

\[ \alpha = 0.5093944013624111 \]

\subsection*{H - What do you have to change in your decision rule for less honest players to be flagged as cheaters?}
The number of wins a player must have out of 100 games to be flagged as a cheater.

\subsection*{I - How does this affect the two types of errors in \\hypothesis testing?}
The first type of error is a \textbf{Type I error}, which is calling someone a cheater who is not cheating.
This will \textit{decrease}, because a higher number of games someone can win before being called a cheater means players can get more lucky without falsely being called a cheater.

\noindent
\\
The second type of error is a \textbf{Type II error}, which is not calling someone a cheater when they are cheating.
This will \textit{increase}, because a higher number of games someone can win before being called a cheater means cheaters can get away with more before being caught.

\end{document}
